\chapter{System Design}
\label{ch:design}

\section{Architectural overview}
Figure~\ref{fig:architecture} shows the two major flows. PDF ingestion creates a local vector index and metadata collection; query handling retrieves context, routes the task, calls the selected agent, and displays the response. The diagram reflects the active imports and calls in \texttt{app.py}, \texttt{src/rag/rag\_pipeline.py}, and \texttt{src/graph/graph\_builder.py}.

\begin{figure}[htbp]
\centering
\resizebox{0.98\textwidth}{!}{\begin{tikzpicture}[
  node distance=8mm and 11mm,
  block/.style={draw,rounded corners,align=center,text width=3.3cm,minimum height=0.8cm,font=\small},
  store/.style={draw,cylinder,shape border rotate=90,aspect=0.25,align=center,text width=2.8cm,minimum height=1.0cm,font=\small},
  flow/.style={-{Stealth[length=2.2mm]},thick},
  node distance=8mm
]
\node[block] (student) {Student};
\node[block,below=of student] (ui) {Streamlit interface\\upload, task, question, answer};
\node[block,below left=13mm and 17mm of ui] (pdf) {Uploaded PDF};
\node[block,below=of pdf] (extract) {Page text extraction\\pypdf};
\node[block,below=of extract] (split) {Word-based splitter\\overlapping chunks};
\node[block,below=of split] (embed) {Sentence Transformer\\normalized vectors};
\node[store,below=of embed] (index) {FAISS L2 index\\and metadata};
\node[block,below right=13mm and 17mm of ui] (query) {Question and selected intent};
\node[block,below=of query] (retrieve) {Hybrid retrieval\\dense + keyword scoring};
\node[block,below=of retrieve] (router) {LangGraph router};
\node[block,below=of router] (agents) {Teacher, quiz, flashcard, summary, planner, doubt agents};
\node[block,below=of agents] (llm) {Groq chat-model wrapper\\when credentials are configured};
\node[block,below=of llm] (response) {Response and source metadata};

\draw[flow] (student) -- (ui);
\draw[flow] (ui) -- (pdf);
\draw[flow] (pdf) -- (extract);
\draw[flow] (extract) -- (split);
\draw[flow] (split) -- (embed);
\draw[flow] (embed) -- (index);
\draw[flow] (ui) -- (query);
\draw[flow] (query) -- (retrieve);
\draw[flow] (index.east) -- ++(8mm,0) |- (retrieve.west);
\draw[flow] (retrieve) -- (router);
\draw[flow] (router) -- (agents);
\draw[flow] (agents) -- (llm);
\draw[flow] (llm) -- (response);
\draw[flow] (response.west) -| (ui.south);
\end{tikzpicture}
}
\caption{Source-backed system architecture of the application}
\label{fig:architecture}
\end{figure}

\section{Component responsibilities}
\begin{table}[htbp]
\centering
\caption{Main runtime components}
\label{tab:components}
\begin{tabularx}{\textwidth}{p{0.27\textwidth}Y}
\toprule
\textbf{Component} & \textbf{Responsibility}\\
\midrule
\texttt{app.py} & Streamlit controls, uploaded-file persistence, task selection, query orchestration, and answer display.\\
\texttt{config.py} & Environment-backed defaults, filesystem paths, and directory creation.\\
\texttt{src/rag} active path & PDF text loading, splitting, embedding, FAISS persistence, hybrid retrieval, and RAG coordination.\\
\texttt{src/graph} & Shared state schema, intent router, agent nodes, and graph compilation.\\
\texttt{src/agents} and \texttt{src/prompts} & Task-specific prompt construction, model invocation, response sanitization, and confidence assignment.\\
\texttt{src/models/llm.py} & Lazy Groq chat-model adapter and user-readable fallback strings.\\
\texttt{src/utils} & Source formatting, prompt instruction, response sanitization, confidence heuristic, PDF extraction, and logging.\\
\bottomrule
\end{tabularx}
\end{table}

\section{Data and control flow}
The compiled graph contains a router node and six terminal agent nodes. The router selects one node from the value returned by \texttt{route\_query}; each task agent consumes the retrieved documents, assembles their text as prompt context, invokes a model wrapper when context is non-empty, and updates the shared state. Figure~\ref{fig:query-workflow} summarizes query-time execution.

\begin{figure}[htbp]
\centering
\resizebox{0.98\textwidth}{!}{\begin{tikzpicture}[
  node distance=8mm and 10mm,
  block/.style={draw,rounded corners,align=center,text width=3.0cm,minimum height=0.85cm,font=\small},
  flow/.style={-{Stealth[length=2.2mm]},thick}
]
\node[block] (input) {User question\\and selected task};
\node[block,right=of input] (retrieve) {RAG retrieval\\top-$k$ context blocks};
\node[block,right=of retrieve] (route) {Intent router\\select one agent};
\node[block,right=of route] (prompt) {Task prompt\\question + context};
\node[block,below=of prompt] (model) {Groq chat model\\or fallback message};
\node[block,left=of model] (clean) {Sanitize response\\compute heuristic confidence};
\node[block,left=of clean] (display) {Streamlit display\\answer + sources + timing};
\node[block,below=of retrieve] (index) {FAISS vectors\\and chunk metadata};

\draw[flow] (input) -- (retrieve);
\draw[flow] (index) -- (retrieve);
\draw[flow] (retrieve) -- (route);
\draw[flow] (route) -- (prompt);
\draw[flow] (prompt) -- (model);
\draw[flow] (model) -- (clean);
\draw[flow] (clean) -- (display);
\end{tikzpicture}
}
\caption{Query-time workflow from user input to displayed answer}
\label{fig:query-workflow}
\end{figure}

\section{State model}
The \texttt{AgentState} TypedDict declares optional fields for the question, retrieved documents, intent, selected agent, response, sources, pages, chat history, confidence, routing label, response time, and retrieved chunk text. At runtime, \texttt{app.py} constructs this state before calling \texttt{graph.invoke}. The agent functions mutate and return the supplied dictionary; no external checkpoint store is configured.

\section{Persistence design}
The configured vector index path is \texttt{vector\_store/faiss.index}; metadata is stored in \texttt{vector\_store/metadata.pkl}. Uploaded PDFs are written to \texttt{data/uploads}. Both paths are ignored by Git, as are logs. The application's conversation history is held in Streamlit session state rather than a durable database.
